\medskip

\textbf{Partie A}

\medskip

\begin{enumerate}
	\item\
	\begin{center}
	\pstree[treemode=R,nodesepA=0pt,nodesepB=2.5pt,treesep = 1cm,levelsep=2.5cm]{\TR{}}
	{\pstree{\TR{$V$~~}\naput{\np{0.6}}}
		{\TR{$L$}\naput{\np{0.2}}
		\TR{$\overline{L}$}\nbput{\np{0.8}}
	}
	\pstree{\TR{$\overline{V}$~~}\nbput{\np{0.4}}}
		{\TR{$L$}\taput{$x$}
		\TR{$\overline{L}$}\nbput{$1-x$}
	}}
	\end{center}
    \item La probabilité qu'un client choisisse un bateau à voile et qu'il ne prenne pas l'option PILOTE est $P\left(V \cap \overline{L}\right)$.
	
	\[P\left(V \cap \overline{L}\right) = P(V) \times P_{V}\left(\overline{L}\right) = \np{0.6} \times \np{0.8} = \boxed{\np{0.48}}\]
	\item Les évènements $V$ et $\overline{V}$ forment une partition de l'univers, donc, d'après la formule des probabilités totales, on a :
	\[P(L) = P\left(V \cap L\right) + P\left(\overline{V} \cap L\right)\]
	On a $P\left(V \cap L\right) = \np{0.6} \times \np{0.2} = \np{0.12}$.
	
	Posons $P\left(\overline{V}\cap L\right) = p$. On a :
	\[P(L) = \np{0.12} + p \iff \np{0.42} = \np{0.12} + p \iff p = \np{0.42}-\np{0.12} = \boxed{\np{0.3}}\]
	On a donc bien $P\left(\overline{V} \cap L\right) = \np{0.3}$.
	\item D'après la formule des probabilités conditionnelles :
	\[P_{\overline{V}}(L) = \dfrac{P\left(\overline{V} \cap L\right)}{P\left(\overline{V}\right)} = \dfrac{\np{0.3}}{\np{0.4}} = \boxed{\dfrac{3}{4} = \np{0.75}}\]
	\item La probabilité qu'un client ait choisi un voilier, sachant qu'il a pris l'option PILOTE est, d'après la formule des probabilités conditionnelles :
	\[P_{L}(V) = \dfrac{P\left(V \cap L\right)}{P(L)} = \dfrac{\np{0.12}}{\np{0.42}} = \boxed{\dfrac{2}{7} \approx \np{0.29}}\]
\end{enumerate}

\medskip

\textbf{Partie B}

\medskip

\begin{enumerate}
	\item \[P\left(L \cap A\right) = P(L) \times P_{L}(A) = \np{0.42} \times \np{0.005} = \boxed{\np{0.0021}}\]
	\smallskip
	
	\[P\left(\overline{L} \cap A\right) = P\left(\overline{L}\right) \times P_{\overline{V}}(A) = (1-\np{0.42}) \times \np{0.12} = \boxed{\np{0.0696}}\]
	\item Les évènements $L$ et $\overline{L}$ forment une partition de l'univers, donc, d'après la formule des probabilités totales, on a :
	\[P(A) = P\left(L \cap A\right) + P\left(\overline{L} \cap A\right) = \np{0.0021} + \np{0.0696} = \np{0.0717}\]
La probabilité que le bateau loué par un client choisi au hasard subisse une avarie est donc \np{0.0717}.
	
Puisque l'entreprise loue \np{1000} bateaux, elle peut s'attendre à 72 avaries.
\end{enumerate}

\medskip

\textbf{Partie C}

\medskip

\begin{enumerate}
	\item Les paramètres de la loi binomiale suivie par $X$ sont $\boxed{n=40}$ et $\boxed{p=\np{0.42}}$.
	\item À l'aide de la calculatrice, on calcule $P(X \geqslant 15)$ :
	\[\boxed{P(X \geqslant 15) \approx \np{0.768}}\]
\end{enumerate}

\bigskip
